\clearpage
\scanpage{11}
\uuline{Typical \unsure{problems|problems} to be solved}

\reversemarginpar
\margin{\tikz[baseline]{\node (a) at (0,0) {\S\,7};\draw (0,0) ellipse (.38 and .59);}}
\fbox{a} R\uline{ings of definition} for structures without \unsure{automorphisms|automorphism}

\margin{\tikz[baseline]{\node at (0,0) {\S\S\,20,21};\draw (0,0) ellipse (.84 and .30);}}
b) Hilbert, Picard without projectivity hypothesis
\normalmarginpar\margin{[Partial results]}\reversemarginpar

\margin{\tikz[baseline]{\node at (0,0) {Chapter \deleted{\unsure{6|V}}$\underline{\overline{\mathrm{VI}}}$};\draw (0,0) ellipse (.81 and .39);}}
Dual of a nonprojective abelian scheme
\normalmarginpar\margin{\fbox{OK. by Raynaud}}\reversemarginpar

\margin{\tikz[baseline]{\node at (0,0) {\S\S\,3,4};\draw (0,0) ellipse (.75 and .42);}}
\fbox{c} Passage to the quotient without flatness hypothesis

\noindent\uuline{\hspace{3cm}}

% slanted
\margin{\tikz[baseline,remember picture]{\node[rotate=13] (mumford) at (0,0) {Mumford};\draw[rotate=13] (0,0) ellipse (.81 and .29);}}
\noindent\(\left[\begin{array}{@{}l@{}}
\tikz[baseline,remember picture]{\coordinate (dlabel);}\text{d) }\tikz[baseline,remember picture]{\coordinate (itemd);}\text{Passage to the quotient by }\mathrm{Gl}\,\unsure{n|m}\text{, acting ``without fixed points''}\\[.5em]
\text{e) Moduli of more or less arbitrary varieties.}\\
\qquad\tikz[baseline,remember picture]{\coordinate (abel);}\text{Case of abelian schemes}
\end{array}\right.\)
\tikz[remember picture,overlay]{\draw[gray,thin](mumford.south) -- (itemd); \draw[gray,thin](mumford.south) to[out=-98,in=178] (abel);}

% slanted
\margin{\tikz[baseline]{\node[align=left,inner sep=1pt] at (0,0) {Mumford -\\Seshadri -\\\unsure{Narashiman|Narasimhan}};\draw (0,0) ellipse (1.01 and .67);}
for vector bundles\thinspace; finite case\thinspace: write up}
\noindent\(\left[\begin{array}{@{}l@{}}
\text{f) Classification of principal bundles }\inserted{\text{over }X/Y}\ \text{(by moduli)}\\
\quad\text{\unsure{with} affine group }\overset{\inserted{G}}{\phantom{G}}\text{ over }Y\text{, one distinguishes the cases}\\
\qquad\tikz[baseline,remember picture]{\coordinate (abelg);}\,G\text{ commutative}\\
\qquad G\text{ not necessarily commutative}
\end{array}\right.\)

\margin{\tikz[baseline,remember picture]{\node (seventeen) at (0,0) {\S\,17};\draw[rotate=14] (0,0) ellipse (.53 and .24);}}
\tikz[baseline,remember picture]{\coordinate (itemg);}g) Special case\thinspace; \(G\) finite\thinspace:

\tikz[remember picture,overlay]{\draw[gray,thin] (seventeen.east) to[out=12,in=182] (abelg); \draw[gray,thin] (seventeen.east) to[out=-21,in=168] (itemg);}
% slanted
\margin{Problem ill-posed\,?}
h) Special case\thinspace: \(X=P^{1}_{Y},\quad G=\mathrm{Gl}(n)\).

\noindent\uuline{\hspace{3.7cm}}

\margin{Chapter $\underline{\overline{\mathrm{VI}}}$\quad\tikz[baseline]{\node at (0,0) {yes};\draw[rotate=13] (0,0) ellipse (.53 and .35);}}
i) Are the formal moduli of an abelian variety \deleted{$A$} \uline{smooth} over the base ring\,?

\noindent\uuline{\hspace{3.8cm}}

\fbox{j)} ``Theorem of the square'' in nonprojective case \(-\, -\, -\, -\)

% margin note written sideways
% Paper watermark through the center.
\margin{\unsure{EXTRA} \unsure{P.J.A|P.J.} \unsure{WOVE}}
% archive sheet number: 10
